\section{Síntesis y ampliación}

\subsection{Puntos clave de esta clase}

\begin{enumerate}
    \item \textbf{Marco unificado SDE}: Los modelos DDPM discretos y SMLD pueden generalizarse naturalmente a SDEs en tiempo continuo. El proceso hacia adelante queda completamente determinado por el coeficiente de deriva $f$ y el coeficiente de difusión $g$.

    \item \textbf{VE-SDE y VP-SDE}:
    \begin{itemize}
        \item VE-SDE (correspondiente a SMLD): $f = 0$, $g = \sqrt{\mathrm{d}[\sigma^2]/\mathrm{d}t}$, con varianza que aumenta monótonamente;
        \item VP-SDE (correspondiente a DDPM): $f = -\frac{1}{2}\beta(t)\mathbf{x}$, $g = \sqrt{\beta(t)}$, con varianza que se mantiene constante.
    \end{itemize}

    \item \textbf{SDE hacia atrás}: Se obtiene directamente del teorema de Anderson; solo es necesario conocer adicionalmente la función score. Una vez entrenada una red score, tanto el proceso hacia adelante como el hacia atrás quedan completamente determinados.

    \item \textbf{Modelos de Difusión Variacional}: A través de la parametrización $(\alpha_t, \sigma_t)$, bajo la condición de que la relación señal-ruido (SNR) disminuya monótonamente, se define una familia infinita de modelos de difusión. VE-SDE y VP-SDE son solo dos instancias de esta familia.

    \item \textbf{EDO de Flujo de Probabilidad}: Cada SDE tiene una EDO correspondiente que genera la misma distribución marginal pero con trayectorias deterministas. Esto proporciona herramientas para el muestreo eficiente y el cálculo de verosimilitud preciso.
\end{enumerate}

\begin{importantbox}{Valor de la perspectiva en tiempo continuo}
El valor del marco SDE en tiempo continuo no radica únicamente en "unificar" los métodos existentes. Lo más importante es que nos ofrece una metodología sistemática para diseñar nuevos modelos de difusión: basta con definir $f$ y $g$ (o equivalentemente, $\alpha_t$ y $\sigma_t$), mientras que lo restante (proceso hacia atrás, EDO, objetivo de entrenamiento) queda determinado automáticamente. Esta flexibilidad tipo "plug and play" ha impulsado una gran cantidad de trabajos posteriores.
\end{importantbox}

\subsection{Direcciones futuras}

El marco en tiempo continuo establecido en esta clase constituye la base teórica para varios temas importantes subsiguientes:

\begin{itemize}
    \item \textbf{Muestreadores eficientes}: Solucionadores numéricos de alto orden basados en la EDO de Flujo de Probabilidad (como DPM-Solver, DEIS, etc.) pueden generar muestras de alta calidad en solo 10--15 pasos;
    \item \textbf{Flow Matching / Rectified Flow}: Un método más conciso para entrenar modelos generativos en tiempo continuo, que puede considerarse una parametrización directa de la EDO de Flujo de Probabilidad;
    \item \textbf{Modelos de Consistencia}: Logran generación en un solo paso mediante la imposición de consistencia entre mapeos en diferentes momentos sobre la trayectoria de la EDO;
    \item \textbf{Evaluación de verosimilitud precisa}: Cálculo del logaritmo de la verosimilitud aprovechando la EDO y el estimador de traza de Hutchinson.
\end{itemize}

\subsection{Lecturas adicionales}

\begin{itemize}
    \item Song, Y., et al. (2021). ``Score-Based Generative Modeling through Stochastic Differential Equations.'' \textit{ICLR 2021.} —Referencia central de esta clase.
    \item Anderson, B. D. O. (1982). ``Reverse-time diffusion equation models.'' \textit{Stochastic Processes and their Applications.} —Teoría clásica de la SDE hacia atrás.
    \item Kingma, D. P., et al. (2021). ``Variational Diffusion Models.'' \textit{NeurIPS 2021.} —Propuesta del marco de Modelos de Difusión Variacional.
    \item Karras, T., et al. (2022). ``Elucidating the Design Space of Diffusion-Based Generative Models.'' \textit{NeurIPS 2022.} —Trabajo importante que explora sistemáticamente el espacio de diseño de modelos generativos basados en difusión en tiempo continuo.
\end{itemize}

%% --- Fin del contenido --- %%

\end{document}
